\section{Challenges and Future Work}\label{sec-challenge}
After proposing the holistic framework of GraphRAG and reviewing it in each domain, we begin in this section by outlining the challenges and opportunities associated with each key component of GraphRAG, including graph construction, retriever, organizer, and generator. Then, we discuss the challenges and opportunities for GraphRAG as a holistic system with its evaluation and application. 

 \subsection{Graph Construction} 

 \begin{itemize}[leftmargin=*]
     \item \textbf{How to construct graphs?} There are numerous ways to construct graphs, yet different tasks or domains may require different graph structures. For example, deciding on the granularity of nodes and edges, as well as which entities or relationships to extract, is critical. This process may also need to address challenges such as entity disambiguation, entity alignment, and coreference resolution. Understanding when a graph structure is necessary, whether to use single or multiple graphs and how to construct them in the most appropriate way for a specific application is essential but often complex. \\
    
    \item \textbf{The format of graphs.} Graphs can be represented in various forms—are these representations equivalent, or do they offer unique advantages? Choosing the most effective representation for a given task can significantly influence performance. \\
    
    \item \textbf{Multi-modal Graphs.} Despite building text-based graphs, the retrieved resources, such as images, audio, or video, can be multi-modal. Constructing a cohesive graph from multi-modal data presents a significant challenge, as it requires integrating diverse data types while preserving meaningful relationships. \\
    
    \item \textbf{Dynamic graphs.} Many real-world scenarios involve dynamic data that evolve over time, which are essential for downstream tasks. Developing strategies to construct, update, and store graphs dynamically while maintaining both efficiency and effectiveness presents a further challenge worthy of exploration.
\end{itemize}
  %  \item \textbf{Challenge 2: Query} 
   % \yu{Deciper the intention}
    %\yu{Entity Ambiguity}

\subsection{Retriever}
\begin{itemize}[leftmargin=*]
    \item \textbf{Differentiating Neural and Symbolic Knowledge}: Graph-structured data usually involves two distinct types of knowledge: symbolic-formatted knowledge, such as relations in knowledge graphs, and neural-formatted knowledge, such as entity names. Developing techniques to differentiate these two types of knowledge and designing corresponding retriever strategies for retrieving these two types of knowledge is worthy of further investigation.

    \item \textbf{Harmonization between internal and external knowledge}: Since GraphRAG is often applied when internal knowledge alone is insufficient to address the task, assessing any overlap between internal and external knowledge using a robust calibration method is crucial. In addition, external knowledge may sometimes conflict with internal knowledge. To handle this, it is essential to design an effective knowledge validation and reconciliation mechanism, allowing selective retrieval and updating of external content as needed.


    \item \textbf{Trade-off among Accuracy, Diversity, and Novelty}: Real-world user retrieval often involves complex intentions. Beyond delivering accurate content to ensure high utility—such as achieving high question-answering accuracy—there may also be a demand for diversity and novelty in the retrieved information. Balancing retrieved content's accuracy, diversity, and novelty remains an open-ended challenge.

    \item \textbf{Reasoning, planning, and thinking along the way}: A real-world retriever may need to dynamically and adaptively update its retrieve process in response to both the initial query and the content retrieved along the way. The question of how to equip the retriever with these adaptive thinking, reasoning, and planning abilities remains a big problem.

\end{itemize}
    
\subsection{Organizer} 
\begin{itemize}[leftmargin=*]
    \item \textbf{Balancing Completeness and Conciseness.} The retrieved graphs may be large, potentially containing significant information that is not related to the query. The Organizer should balance the need for complete information with the risk of overwhelming the model. This requires techniques for pruning irrelevant nodes and edges while retaining essential context, which is especially challenging in large or noisy graphs. Additionally, some knowledge may already be captured by LLMs, raising the question of how to identify and remove redundant information to further improve efficiency. \\ 
    
    \item \textbf{Optimal Data Structuring:} Determining the most effective way to structure the retrieved content is a challenge. For example, deciding how to convert a structured graph into a format that the generator can leverage, how to arrange the order of retrieved content, and how to preserve the original data structure for structure-sensitive tasks are all important considerations. Different tasks may benefit from different structuring methods. \\
    
    \item \textbf{Aligning Different Resources:} Retrieved content may come from various sources and in diverse formats, such as text, graphs, and images. Aligning these components effectively to help the generator poses a significant challenge, especially when integrating multiple data modalities. \\
    
    \item \textbf{Data Augmentation: } Incorporating data augmentation within the Organizer involves enriching the retrieved graph content with nodes, edges, or features to improve the robustness of the model and improve downstream tasks. However, this process requires balancing real and augmented data to avoid introducing irrelevant or redundant information.
\end{itemize}  
    
\subsection{Generator}
\begin{itemize}[leftmargin=*]
    \item \textbf{Correct Format for Prompting.} The content retrieved after organization varies significantly in format — such as texts, triplets, or graphs — while current LLMs can only process text inputs. Exploring the most effective format for optimal LLM performance on specific tasks and designing more flexible generators that go beyond text input could be worthwhile avenues for research.

    \item \textbf{Structural Encoding.} When the retrieved content is a subgraph and the downstream generator is an LLM, ensuring that the LLM can interpret the graph's structural information is essential. Designing effective structural encodings and integrating them into token embeddings pose a key challenge. Although some previous works have incorporated various graph encodings into the text decoding process, no systematic study has yet demonstrated whether LLMs can distinguish between these encodings and accurately recognize their corresponding geometric structures.
\end{itemize}  


\subsection{GraphRAG as a System} Rather than individual components, GraphRAG is a system. Designing an efficient and cohesive GraphRAG system presents additional challenges: \\
\begin{itemize}[leftmargin=*]
    \item \textbf{Integration Across Components:} Ensuring seamless interaction among the Graph Construction, Retriever, Organizer, and Generator components is essential. Each part must operate harmoniously to maintain efficiency and accuracy, which can be challenging to optimize the system globally. \\
    \item \textbf{Scalability:} As data volumes and query demands increase, each component of the GraphRAG system, such as Graph Construction, Retriever, Organizer, and Generator, must efficiently handle larger, more complex graphs without compromising performance.For example, efficient graph storage, optimized querying (e.g., subgraph sampling and pathfinding), streamlined organization of retrieved components, and responsive generation are all essential. Additionally, training, serving, fine-tuning, and evaluating within GraphRAG demand sophisticated engineering on modern hardware and software stacks, including efficient training, data-efficient fine-tuning, communication-efficient algorithms, implementation of reinforcement learning with human feedback, GPU acceleration and other specialized hardware, model compression for deployment, and online maintenance. \\
   % \textbf{3) Security:} As a system-level concern, security involves ensuring that GraphRAG is protected from potential misuse.  Key security challenges include safeguarding against attempts to recover sensitive training data embedded within models and defending against adversarial inputs designed to manipulate models into producing target outputs. Developing robust defenses against these vulnerabilities is essential for maintaining the integrity and confidentiality of the GraphRAG system. \\
   % \textbf{4) Trustworthiness and Explainability:} For users to trust GraphRAG outputs, the system should provide insights into how answers were derived. This is especially challenging in complex graph-based retrieval and generation systems, where multiple relationships and entities contribute to a response. Ensuring transparency and explainability is crucial, particularly for applications in sensitive domains like healthcare, finance, and law. \hy{May move this part to the Trustworthiness}
    \item \textbf{Trustworthiness:}
    Real-world GraphRAG systems often operate in high-stakes domains such as education~\cite{Thway2023Battling}, healthcare~\cite{Xiong2023Benchmarking, Zakka2023Almanac}, and law~\cite{Wiratunga2023CBR}, which imposes multiple desires when deploying GrapRAG systems. Therefore, deeply understanding the broader desires beyond merely optimizing utility is essential to ensure effective deployment. As motivated by trustworthy and safety research in other fields~\citep{chen2024fairness, huang2024position, wang2024safety}, these key goals beyond utility include reliability, robustness, fairness, and privacy. Although previous works have initiated the exploration of this multi-purpose optimization of RAG~\cite{wagle2023empirical, zhou2024trustworthiness}, they mostly focus on RAG without dedicated investigation of how the relational information captured by GraphRAG could cause additional trustworthy and safety concerns. 
    \begin{itemize}
        \item \textbf{Reliability}~\cite{kumar2023conformalpredictionlargelanguage, quach2024conformallanguagemodeling, su2024apienoughconformalprediction, ye2024benchmarkingllmsuncertaintyquantification, li2023traq, ni2024towards}: Reliability requires the system to deliver consistently low error rates across different scenarios. In RAGs, uncertainty quantification presents two primary challenges: Firstly, during the generation phase, uncertainty stems from the inherent probabilistic generation of LLMs. Standard techniques such as conformal prediction have been applied here with few adaptations~\cite{kumar2023conformalpredictionlargelanguage, quach2024conformallanguagemodeling, su2024apienoughconformalprediction, ye2024benchmarkingllmsuncertaintyquantification}. The general idea is to estimate the non-conformal score based on the calibration set and filter out those high-risk answers during the testing phase. Secondly, the retriever and its interaction with LLMs (e.g., treating LLMs as the reasoning agent to perform adaptive retrieval) also possess uncertainty~\cite{li2023traq, ni2024towards}. Unlike RAGs, retrievers and generators in GraphRAG demonstrate the multi-hop nature, which raises new requirements for multi-stage uncertainty quantification. For example, the error rate of one-step generation might accumulate after multi-step generation, and how to calibrate this accumulated error at the global level is the key challenge. The very first work~\cite{ni2024towards} proposed the learn-then-test framework to guarantee the global error rate. However, this solution may not consider the uncertainty of human-LLM interactions. Moreover, the additional learn-then-test component takes additional time and computational load for calibration. Future research aims to develop methods that can quantify and calibrate the uncertainty of (Graph)RAG as a system holistically.

        \item \textbf{Robustness}~\cite{fang2024enhancing, xiang2024certifiably, xu2024unveil, yoran2023making, dziri2024faith, zhang2024darg, zhu2024dyval}: Robustness aims the system to deliver equal-quality response under extreme scenarios such as the presence of noisy and irrelevant content. Existing research on RAGs has highlighted the significant degradation in LLMs' performance when retrieved content includes noise or irrelevant information~\cite{yoran2023making, fang2024enhancing}. Some studies address this challenge by adversarially training LLMs in noisy environments. However, there has been limited exploration of LLMs' robustness from the structural perspective. For example, would the LLM reasoning performance change when the underlying reasoning graph gets perturbed~\cite{dziri2024faith, zhang2024darg, zhu2024dyval}? Addressing these issues requires a deeper focus on the interplay between structural integrity and the robustness of GraphRAG systems, paving the way for robust performance in tasks requiring complex reasoning and planning.
        
        
        % the out-of-distribution not only happens at the instance level but also at the relation level. For example, \citet{zhang2024darg} has shown that the mathematical reasoning capability of LLMs gradually decreases as the reasoning graph structure (e.g., graph depth and width) becomes more complex. Beyond generalizability as one aspect of robustness, this relational information provides malicious actors additional channels to design adversarial attacks, such as injecting malicious passage nodes to dismantle graph-structured data sources and misguide graph traversal. 

        \item \textbf{Safety}~\cite{zou2024poisonedragknowledgecorruptionattacks, xue2024badrag, wang2024poisonedlangchainjailbreakllms, deng2024pandorajailbreakgptsretrieval, xiang2024certifiably}: Recent studies have highlighted that LLMs are susceptible to various adversarial attacks. Techniques such as prompt manipulation, hint injection, and input perturbation enable attackers to bypass safety mechanisms and exploit vulnerabilities, posing substantial social risks. RAGs that combine the power of LLMs with external databases introduce unique safety challenges. Existing works in RAGs have developed several threats (e.g., adversarial passage injection~\cite{zou2024poisonedragknowledgecorruptionattacks}, group-query targeting~\cite{xue2024badrag}, and jailbreak attacks~\cite{wang2024poisonedlangchainjailbreakllms, deng2024pandorajailbreakgptsretrieval}) and defense strategies (such as isolate-then-aggregate strategy~\cite{xiang2024certifiably}). However, all of them have overlooked the structural vulnerabilities inherent in graph-structured data. For instance, attackers can exploit the network structure of graph-based databases by leveraging principles from network science to design highly impactful attacks. Attackers can maximize their influence on the system's behavior by strategically targeting nodes with specific structural properties, such as high degree or centrality. This highlights the need for robust defenses that account for both content-based and structural threats in RAG systems. 
        
        
        % How to ensure the quality of retrieved content under this unexpected scenario remains unexplored. Despite previous works, they mostly focus on the adversarial robustness of the generator.  Although previous works have developed multiple Fair RAGs~\cite{kong2024mitigating, xue2024badrag, shrestha2024fairrag}, they are mainly for image, and the fairness is mainly through the sensitive attributes. However, fairness issues could also be sourced from structural and generated at different perspectives, such as group and individual fairness~\cite{dong2023fairness, wang2022improving}; how to guarantee fairness at structure in different levels and averaged in RAG remains never explored. As a system-level concern, security involves ensuring that GraphRAG is protected from potential misuse.  Key security challenges include safeguarding against attempts to recover sensitive training data embedded within models and defending against adversarial inputs designed to manipulate models into producing target outputs. Developing robust defenses against these vulnerabilities is essential for maintaining the integrity and confidentiality of the GraphRAG system. For users to trust GraphRAG outputs, the system should provide insights into how answers were derived. This is especially challenging in complex graph-based retrieval and generation systems, where multiple relationships and entities contribute to a response. Ensuring transparency and explainability is crucial, particularly for applications in sensitive domains like healthcare, finance, and law. Addressing these complexities requires careful consideration of the structure's and trustworthiness interplay, which benefits many real-world applications, especially for high-risk scenarios.

        \item \textbf{Privacy}~\citep{chen2024survey, zafar2023building, qian2017social}: Privacy is a critical concern in RAG systems~\citep{zeng2024good, zhou2024trustworthiness}, especially when operating in domains involving sensitive or personal data, such as healthcare, education, or finance~\citep{yao2024survey, das2024security, murali2023towards}. Unlike traditional RAG systems, GraphRAG introduces additional privacy risks due to the relational nature of graphs, which may inadvertently reveal private information through connections or patterns. For instance, even if a sensitive node is protected, it may still be retrieved indirectly via its neighbors. Additionally, in graphs exhibiting homophily—where connected nodes tend to share similar attributes—sensitive information can be inferred from neighboring nodes. The use of GNNs for graph encoding further exacerbates these risks. The message-passing mechanism in GNNs propagates features across nodes and edges, potentially leading to the leakage of sensitive or confidential information during the propagation process~\citep{zhang2024survey}.  Addressing these challenges requires advanced privacy-preserving techniques that account for the interconnected nature of graphs, rather than treating the data as independent and identically distributed (iid).  Privacy protection must be designed across the entire GraphRAG pipeline, from graph construction to generation, ensuring that sensitive information remains secure while preserving the graph's utility for downstream tasks.


        \item \textbf{Explainability}: Explainability is essential for fostering trust in RAG systems, especially in high-stakes domains like law, healthcare, and finance~\citep{zhao2024explainability, cambria2024xai, yu2023temporal, elsborg2023using}. Compared to traditional RAG, GraphRAG offers enhanced explainability through the explicit relationships encoded between nodes in the graph~\citep{abu2024supporting}. These explicit connections allow the system to generate clear and interpretable explanations tailored to specific tasks, making it more transparent and trustworthy for end users. For example, in the mutli-hop question answering, the GraphRAG can be leveraged to generate reasoning paths to the question~\citep{luo2023reasoning, keqing, sunthink, jiang2024hykge}. These paths provide step-by-step explanations, allowing users to understand how intermediate conclusions were derived and verify the relevance and correctness of the system’s logic. Similarly, in molecular property prediction, specific subgraphs representing functional groups or structural motifs can be utilized to explain predictions, linking molecular features to observed properties or behaviors. However, achieving such explainability requires GraphRAG to retrieve reasonable and contextually relevant subgraphs, which remains a challenge. Ensuring that explanations remain faithful to the underlying model logic while balancing comprehensiveness and simplicity will be critical challenges for future research. 

        % \item \textbf{Accountability}
    \end{itemize}
\end{itemize}

\subsection{ Evaluation of GraphRAG.} Evaluating the performance of a GraphRAG system is challenging due to its complex, multi-component nature. Standard benchmarks may not fully capture the nuances of graph-based construction, retrieval, organization, and generation, so tailored benchmarks are essential to assess each component and their overall impact on the system. \\ 

\begin{itemize}[leftmargin=*]
    \item \textbf{Component-Level Optimality:} Each component’s performance directly influences the GraphRAG system as a whole. Evaluating each component, such as Graph Construction, Retriever, Organizer, and Generator, requires specific designs suited to their unique roles. This may involve constructing different datasets and evaluation metrics that align with the intended function of each component.
    
    \item \textbf{End-to-End Benchmarks:}  To assess the system’s overall effectiveness, comprehensive end-to-end benchmarks are essential. These should evaluate the quality of generated outputs, system response time, efficiency, and resource utilization, providing a holistic view of GraphRAG’s performance in real-world applications.
    
    \item \textbf{Task and Domain-Specific Evaluation:} Different tasks and domains may impose unique requirements on the GraphRAG system, necessitating specialized designs for each component. Methods that perform well in one domain may not generalize effectively to others, highlighting the need for diverse benchmarks. Additionally, Multi-task and multi-domain evaluations help determine the system’s adaptability and effectiveness across varied contexts.
    
    \item \textbf{Trustworthiness benchmark.} Ensuring the trustworthiness of the GraphRAG system is critical, especially in applications where decisions rely heavily on accurate and unbiased information. A trustworthiness benchmark should evaluate aspects such as robustness against adversarial inputs, data privacy protection, and transparency in how answers are derived, ensuring that outputs are explainable and reliable.
\end{itemize}     


\subsection{New Applications} While we have introduced applications of GraphRAG in several domains and tasks, there are still many domains that can leverage GraphRAG, such as Code generation~\cite{liu2024codexgraph} and robust cyber defense~\cite{rahman2024retrieval}. Extending its use to other domains is promising but presents unique challenges. Determining how to adapt GraphRAG effectively for new fields requires understanding the specific requirements, data structures, and graph configurations unique to each application area. Developing tailored strategies to construct, retrieve, organize, and generate data for diverse domains is essential for maximizing GraphRAG’s adaptability and effectiveness.

     
     %Below we give some examples for potenital new applications. 
%\begin{itemize}
 %   \item \textbf{Graph RAG as a Memory Management Tool in Multi-Agent Systems.} In multi-agent systems, effective memory management~\cite{han2024llm, li2023mot} is crucial for ensuring that agents can collaboratively interact, learn from past experiences, and make informed decisions in real time. 
 %   Traditionally, memory systems in multi-agent frameworks are challenged by the need to represent vast amounts of memory, handle dynamic memory updates, and ensure efficient access to relevant memory on demand. To address these challenges, Graph RAG (Graph Retrieval-Augmented Generation) offers a powerful solution by combining the strengths of graph-structured data representation with the retrieval-enhanced capabilities of modern generation models. 
    
  %  \textbf{1) Representation of vast amounts of memory}: In multi-agent systems, agents generate and share extensive memories, encompassing individual observations, interactions, and experiences over time. Traditional approaches struggle to represent this vast information in a structured, accessible way. The challenge lies in organizing this data to facilitate both individual and collective recall. Knowledge graph in Graph RAG addresses this by organizing memory as a graph structure, where memories are represented as nodes and their relationships as edges. This enables a structured and inherently scalable approach to representing memory in a relational, interpretable format.
    
 %   \textbf{2) Dynamic memory updates}: As agents interact with dynamic environments, their memory requirements continuously change. This necessitates a flexible memory system that allows for the addition and deletion of memories in real time. Agents need to be able to learn from new information while discarding outdated or irrelevant data to avoid memory overload. Graph RAG provides a solution by enabling real-time updates to the memory graph, allowing for the insertion of new nodes and edges to represent fresh knowledge and the removal of obsolete memory. This adaptability ensures that memory remains relevant and manageable as agents evolve.

 %   \textbf{3) Efficient access to relevant memory}: Efficient retrieval is essential in memory management, particularly when agents must quickly access contextually relevant information to make timely decisions. In multi-agent systems, agents often need to retrieve specific memory items based on complex contextual cues, which can be computationally intensive. Graph RAG leverages retrieval-augmented generation capabilities to provide fast, context-driven access to memory. By identifying and prioritizing the most relevant nodes and edges within the graph, agents can efficiently recall essential information on demand, optimizing response times and decision quality.
    
    
    % By leveraging a graph-based structure, Graph RAG allows for a natural representation of memory as interconnected nodes and edges, where each node can encapsulate distinct memories, experiences, or knowledge components, and edges represent the relationships between them. This setup not only mirrors how human memory networks work but also enables highly efficient memory storage and retrieval processes, essential for complex, interactive tasks. 

 %   \item \textbf{ Knowledge Verification for LLMs Using Graph RAG.} The generated content from Large language models (LLMs) are not evidence-based, leading to a risk of hallucinations~\cite{huang2023survey}-generated information that may be inaccurate or entirely fabricated. This issue can be particularly concerning in critical fields, such as healthcare~\cite{ahmad2023creating}, where erroneous information could have serious consequences. To address this challenge, Graph RAG can be deployed as a knowledge verification tool. By constructing a domain-specific knowledge graph (KG)-for example, in the medical field-Graph RAG enables retrieval of evidence-based, verified content that can support or counter-check LLM-generated responses. When a query related to the domain is made, relevant nodes and edges from the KG can be retrieved, providing robust, factual support for the response. The advantage of Graph RAG lies in its cross-document retrieval and reasoning capabilities. Unlike simple retrieval from a static database or individual document, a graph structure allows the system to draw connections across multiple sources, integrating knowledge in a cohesive and comprehensive way. By incorporating Graph RAG for knowledge verification, LLMs can provide responses grounded in evidence-based information. This approach has the potential to improve accuracy and reliability in high-stakes domains, minimizing the risk of critical hallucinations and reinforcing trust in LLM-generated content.

%    \item \textbf{Harmonizing Symbolic and Neural Knowledge} 
 %   GraphRAG integrates two distinct types of knowledge: symbolic-formatted knowledge, such as knowledge graphs and semi-structured knowledge bases~\cite{wu2024stark}, and neural-formatted knowledge, stored as the parameters of LLMs. A deeper understanding of these two knowledge formats—how they store, manage, and interact with information—can result in better harmonization between these two knowledge. From an understanding perspective, it is valuable to explore questions such as: Given a knowledge graph, how much of its information is already represented within an LLM? Conversely, how much of an LLM's knowledge is captured in a knowledge graph? The insights derived by answering these knowledge-checking questions can aid in refining knowledge management approaches, including selectively pruning knowledge graphs to remove redundant information and selectively fine-tuning LLMs with knowledge that fills gaps in their understanding, which has the significant potential to improve the efficiency and performance of GraphRAG systems. Furthermore, from the model-centric perspective, it can inform us on how to design better GraphRAG systems. For example, having prior knowledge about which region LLM-powered generator knows less about the graph could help us design adaptive graph traversal that selectively visits neighbors and retrieves their contents. For certain neighbors that have already been well-formed by LLMs, we avoid the traversal to save more money, while for other unknown areas, we allow more money to be visited. Although existing works have started investigating knowledge checking, they all focus on the neural-formatted knowledge and general RAG, with no exploration of symbolic knowledge and GraphRAG.
    
    





\section{Conclusion}\label{sec-conclusion}

In this survey, we first introduced the demand and rationale behind GraphRAG, highlighting its ability to enhance retrieval-augmented generation by integrating graph-structured information. We then unified the architectural designs of existing GraphRAG approaches into a holistic framework, comprising five key components: graph construction, retriever, organizer, generator, and data source. For each component, we reviewed representative techniques. Given the diversity of graph structures and their applications across different domains, we also explored GraphRAG designs tailored to specific domains. By reviewing GraphRAG applications across varied domains—from knowledge graphs and document graphs to scientific and social graphs, we illustrated how its flexibility allows it to meet unique demands and address a wide range of tasks. Finally we discuss challenges and opportunities that have the potential to push the boundaries for GraphRAG. 

%However, implementing an efficient and effective GraphRAG system presents unique challenges, as discussed across its different components: graph construction, retrieval optimization, organization, and generation. Each component introduces distinct demands, and the system’s overall effectiveness relies on seamless integration and adaptability to diverse data types, domains, and tasks. We outlined critical evaluation considerations, including component-level and end-to-end benchmarks, as well as domain-specific assessments to ensure performance across various contexts. Additionally, we discussed challenges related to trustworthiness and the need for expanding GraphRAG’s applications, highlighting areas for future research and development.

%The potential applications of GraphRAG span multiple domains, from question answering and recommendation systems to complex planning and reasoning tasks and knowledge-based fields like healthcare and finance. As this technology evolves, developing domain-adaptive and efficient GraphRAG systems will be key to expanding its use in real-world applications. Future research may focus on overcoming the highlighted challenges, refining evaluation frameworks, and enhancing scalability, efficiency, and trustworthiness. By addressing these areas, GraphRAG can further improve information retrieval and generation, providing a valuable tool for advancing AI capabilities across diverse fields.